% Unicode corpus (M3): xeCJK with Japanese, Harano Aji by file name.
\documentclass{article}
\usepackage{xeCJK}
\setCJKmainfont[BoldFont=HaranoAjiMincho-Bold.otf]{HaranoAjiMincho-Regular.otf}
\setCJKsansfont[BoldFont=HaranoAjiGothic-Bold.otf]{HaranoAjiGothic-Regular.otf}
\begin{document}
日本語の文章です。ひらがな、カタカナ、漢字が混ざっています。「括弧」と句読点。
\textsf{ゴシック体}と\textbf{太字}。English words in between 日本語の文。
この段落は長くして、日本語の行分割と約物の処理を確かめるためのものです。
\end{document}
